\chapter{核衰变和核反应}

\problem{7.1}{半衰期为 $30.2\ \mathrm a$ 的 $1\ \mathrm{mg}$ $^{137}\mathrm{Cs}$ 的放射性活度是多少？每秒放出多少 $\beta$ 射线和 $\gamma$ 射线？}
% layout
\answerarea

\problem{7.2}{$1\ \mathrm s$ 内测量到 $^{60}\mathrm{Co}$ 放射源发出的 $\gamma$ 射线有 3700 个，设测量效率为 $10\%$，求它的放射性活度。已知它的半衰期为 $5.27\ \mathrm a$，求它的质量。}
% layout
\answerarea
\newpage

\problem{7.3}{中午 12:00 时试管中 \nucl{25}{11}{Na} 核（$\beta$ 放射性，$T_{1/2}=60\ \mathrm s$）是 $1\ \mu\mathrm g$，求试管内原有的钠原子数到 12:10 还有多少？估算一个钠原子衰变的时间。}
% layout
\answerarea

\problem{7.4}{从 $1.10\ \mathrm{kg}$ 铀矿中含有 $0.20\ \mathrm{kg}$ 铅（$^{206}\mathrm{Pb}$）这一事实估计地球的年龄。已知天然铀含 $99.28\%$ $^{238}\mathrm U$、$0.714\%$ $^{235}\mathrm U$ 和 $0.006\%$ $^{234}\mathrm U$，$^{206}\mathrm{Pb}$ 是 $^{238}\mathrm U$ 衰变系列的稳定产物，$^{238}\mathrm U$ 的半衰期为 $4.5\times10^9\ \mathrm a$，衰变系列中其他核素的半衰期均比它小很多。}
% layout
\answerarea
\newpage

\problem{7.5}{在用 $^{14}\mathrm C$ 测量一具古尸的工作中，已知 $^{14}\mathrm C$ 的 $T_{1/2}=5730\ \mathrm a$，$^{14}\mathrm C$ 的浓度为 $10^{-12}$。
\begin{enumerate}
\item 若取样 $50\ \mathrm g$ 碳测量 $1\ \mathrm h$，得到 $\beta$ 放射性计数为 2500 个，不考虑本底计数，设探测效率为 $90\%$，求它的死亡年代。
\item 若取样 $50\ \mathrm g$ 碳，要求测量精度为 $\pm50\ \mathrm a$，需要多长测量时间？
\item 测量 $1\ \mathrm h$，精度要求同上，问取碳样多少？
\item 设 $N_c$、$N_t$ 和 $N_b$ 分别为真实计数、测量到的总计数和本底计数，即有 $N_c=N_t-N_b$，$N_c$ 的统计误差为
\[
\Delta N_c=\sqrt{(\Delta N_t)^2+(\Delta N_b)^2}=\sqrt{N_t+N_b}.
\]
如果本底计数是 $4000\ \mathrm{h^{-1}}$，测量 $1\ \mathrm h$，精度要求仍同上，要求取碳样多少？
\end{enumerate}}
% layout
\answerarea
\newpage

\problem{7.6}{厚度为 $0.2\ \mathrm{cm}$、面积为 $1\ \mathrm{cm^2}$ 的锡片，放在反应堆热中子道中照射 $1000\ \mathrm h$，以生产放射性核素 \nucl{113}{50}{Sn}。已知 $^{112}_{50}\mathrm{Sn}(\mathrm n,\gamma){}^{113}_{50}\mathrm{Sn}$ 反应截面 $\sigma_\gamma=0.8\ \mathrm b$，$^{112}\mathrm{Sn}$ 核素丰度为 $0.96\%$，锡的密度 $\rho=7.29\times10^3\ \mathrm{kg\,m^{-3}}$，垂直投射的热中子通量为 $5\times10^{16}\ \mathrm{m^{-2}\,s^{-1}}$，$^{113}\mathrm{Sn}$ 的衰变纲图如图所示，\nucl{113}{49}{In} 的 $(1/2)^-$ 能级的内转换系数 $\alpha=0.43$。
\begin{enumerate}
\item 求出堆时 $^{113}\mathrm{Sn}$ 的活度 $A$；
\item 求 5 天以后 $\gamma$ 射线（$393\ \mathrm{keV}$）的强度；
\item 第 5 天从母体 Sn 中分离出子体 $^{113}\mathrm{In}$ 核素，制成 100 份，每份注入一个人体脏器，以诊断疾病，求经一天后体内剩余的 $^{113}\mathrm{In}$ 放射性核素，不考虑人体新陈代谢效应；
\item 求发射的 $\gamma$ 射线的多极性。
\end{enumerate}
\fig{7.6}}
% layout
\answerarea
\newpage

\problem{7.7}{$^{239}\mathrm{Pu}$ 的半衰期可用下法测定：将 $m=120.1\ \mathrm g$ 的 $^{239}\mathrm{Pu}$ 球置于液氦中，液氦体积足以阻止 $\alpha$ 粒子。测得液氦蒸发功率为 $0.231\ \mathrm W$，已知 $E_\alpha=5.144\ \mathrm{MeV}$，求 $^{239}\mathrm{Pu}$ 的半衰期。}
% layout
\answerarea

\problem{7.8}{放射性核素 $^{238}_{94}\mathrm{Pu}$ 已作为宇宙飞行用能源，它是一种半衰期为 $90\ \mathrm a$ 的 $\alpha$ 放射性，问：
\begin{enumerate}
\item 为什么 $^{238}\mathrm{Pu}$ 放射 $\alpha$ 射线而不是氘核？
\item $\alpha$ 射线的能量约为 $5.5\ \mathrm{MeV}$，如有 $238\ \mathrm g$ $^{238}\mathrm{Pu}$，它们释放的功率是多少？
\item 上述能量是仪器运行能量的 8 倍，能用多长时间？
\end{enumerate}}
% layout
\answerarea
\newpage

\problem{7.9}{$^{238}_{94}\mathrm{Pu}$ 主要有三种能量的 $\alpha$ 射线，它们是 $5.499\ \mathrm{MeV}$、$5.457\ \mathrm{MeV}$ 和 $5.358\ \mathrm{MeV}$，分别跃迁到 $^{234}_{92}\mathrm U$ 的基态和第一、第二激发态。求：
\begin{enumerate}
\item $^{234}\mathrm U$ 的第一、第二激发能级的能量；
\item $^{238}\mathrm{Pu}$ 基态和 $^{234}\mathrm U$ 的三个能级的自旋和宇称；
\item 产生的 $\gamma$ 射线能量和多极性。
\end{enumerate}}
% layout
\answerarea

\problem{7.10}{在一个为验证大统一理论的质子衰变实验中，使用了一个含有 $10000\ \mathrm t$ 纯水的水池。
\begin{enumerate}
\item 用切伦科夫辐射探测，设效率为 $100\%$，质子寿命为 $10^{32}\ \mathrm a$，束缚在核内的质子以与自由质子一样的速率衰变，试估计一年内可观测到多少衰变；
\item 设有一种衰变方式为 $\mathrm p\to\pi^0+\mathrm e^+$，$\pi^0\to\gamma+\gamma$，计算光子的最大和最小能量。
\end{enumerate}}
% layout
\answerarea
\newpage

\problem{7.11}{确定核素 \nucl{80}{35}{Br} 是否可能发生 $\beta^-$、$\beta^+$ 衰变和轨道电子俘获，是否可能有 $\alpha$ 衰变和发射中子的衰变，并计算各衰变能。}
% layout
\answerarea

\problem{7.12}{$^{82}\mathrm{Se}$ 可经双 $\beta$ 衰变成 $^{82}\mathrm{Kr}$，已知 $M(^{82}\mathrm{Se})=81.916639\ \mathrm u$、$M(^{82}\mathrm{Kr})=81.91344\ \mathrm u$，求衰变能。}
% layout
\answerarea
\newpage

\problem{7.13}{\begin{enumerate}
\item $^8\mathrm B$ 可经 $\beta^+$ 衰变到 $^8\mathrm{Be}$，如图所示。这是太阳产生中微子的一个来源，简要说明哪些种 $\beta^+$ 衰变构成中微子谱的主要来源。
\item 除上述外，太阳还通过 $\mathrm p+\mathrm p\to\mathrm d+\beta^++\nu_e$ 聚变反应产生中微子，试问这些中微子作用在 $^1\mathrm H$ 靶和 $^2\mathrm H$ 靶上能否产生反应？能产生什么反应？已知下列原子的质量剩余（以 MeV 作单位）为
\[
\begin{array}{ccccc}
{}^1\mathrm H&\mathrm n&{}^2\mathrm H&{}^8\mathrm B&{}^8\mathrm{Be}\\
7.289&8.071&13.136&22.922&4.942
\end{array}
\]
\end{enumerate}
\fig{7.13}}
% layout
\answerarea
\newpage

\problem{7.14}{原子核 $^{69}\mathrm{Zn}$ 处于能量为 $436\ \mathrm{keV}$ 的同质异能态时，试求它发射 $\gamma$ 光子后的反冲动能和发射内转换电子后的反冲动能；若 $^{69}\mathrm{Zn}$ 处于高激发态，可以发射中子，试求发射 $436\ \mathrm{keV}$ 中子后的反冲能。}
% layout
\answerarea

\problem{7.15}{$^{182}_{73}\mathrm{Ta}$ 的激发态能量为 $137\ \mathrm{keV}$，寿命为 $2\times10^{-15}\ \mathrm s$，用同样能量的 $\gamma$ 射线入射到 $^{182}\mathrm{Ta}$ 核时，能否使核共振激发？}
% layout
\answerarea
\newpage

\problem{7.16}{反应 $^{10}\mathrm B+\mathrm d\to{}^8\mathrm{Be}+\alpha+17.85\ \mathrm{MeV}$，当 d 束能量为 $0.6\ \mathrm{MeV}$ 时，在 $90^\circ$ 方向观测到四种能量的 $\alpha$ 射线，它们是 $12.2$、$10.2$、$9.0$ 和 $7.5\ \mathrm{MeV}$，求 $^8\mathrm{Be}$ 的激发能。}
% layout
\answerarea

\problem{7.17}{氘核吸收 $6\ \mathrm{MeV}$ $\gamma$ 射线后发生分解，如果在与入射光子成 $90^\circ$ 角方向上探测到质子，求反应 $Q$ 值和氘核结合能、中子和质子的动能以及中子的飞出方向。}
% layout
\answerarea
\newpage

\problem{7.18}{已观测到下列两核反应：
\[
{}^{14}\mathrm N+\alpha={} ^{17}\mathrm O+\mathrm p-1.26\ \mathrm{MeV},
\]
\[
{}^{16}\mathrm O+\mathrm d={} ^{14}\mathrm N+\alpha+3.13\ \mathrm{MeV}.
\]
如用 d 打 $^{16}\mathrm O$ 时也形成 $^{17}\mathrm O+\mathrm p$，计算 d 能量为 $1.00\ \mathrm{MeV}$ 时在实验室系 $0^\circ$ 和 $90^\circ$ 方向射出的质子的动能。}
% layout
\answerarea

\problem{7.19}{在什么条件下反应 $^7\mathrm{Li}(\mathrm p,\mathrm n){}^7\mathrm{Be}$ 能实现？在这个阈值下的中子实验室能量是多少？}
% layout
\answerarea
\newpage

\problem{7.20}{探测器面积为 $2\ \mathrm{cm^2}$，覆盖一层自发裂变重核，平均寿命为 $\frac13\times10^9\ \mathrm a$，每秒观测到 20 个裂变事例。若放到通量为 $10^{11}\ \mathrm{cm^{-2}\,s^{-1}}$ 的均匀中子流中，探测到的裂变数为 $120\ \mathrm{s^{-1}}$，试问中子诱发裂变截面为多大？}
% layout
\answerarea

\newpage

\problem{7.21}{1987 年 3 月在两个地下大中微子探测器中同时记录到一个由若干中微子事例组成的爆发，它们出现在几秒钟的隔内，这是由于一个靠近银河系边缘的超新星 1987A 的突然爆发所发出的 $10\ \mathrm{MeV}$ 能区的 $\nu$ 和 $\bar\nu$ 引起的。各个中微子探测器是由一个在约 $1000\ \mathrm m$ 深的矿井中的盛有约 $5000\ \mathrm t$ 纯水的容器组成的，水用一大组光电倍增管围绕，以探测相对论带电粒子产生的切伦科夫辐射。$10\ \mathrm{MeV}$ 能区的中微子与纯水主要作用是
\[
\text{（a）}\ \nu+\mathrm e\to\nu+\mathrm e,\qquad
\text{（b）}\ \bar\nu+\mathrm p\to\mathrm n+\mathrm e^+.
\]
\begin{enumerate}
\item $10\ \mathrm{MeV}$ 的中微子在（a）作用中产生的出射电子最大能量是多少？
\item 两反应中的出射带电粒子在质心系中有各向同性分布，哪种反应的出射电子在实验室系中角分布将指示入射中微子方向？
\item 两探测器的事例爆发被推断为 $\nu$（实际是 $\bar\nu$），它们的能量为 $10\!\sim40\ \mathrm{MeV}$。设所有这些中微子是同一时刻从那个超新星爆发中发射的，在到达地球时通过约 $170000$ 光年的距离，它们到达时间有 $2\ \mathrm s$ 弥散，试估计中微子质量上限。
\end{enumerate}}
% layout
\answerarea
\newpage

